% 分野別類題: 連立線形微分方程式 解答例
% コンパイル: TEXINPUTS="../../../00_common:$TEXINPUTS" latexmk -xelatex answers.tex
\documentclass[a4paper,11pt]{article}
\usepackage{preamble}
\usepackage{shoakubox}

\begin{document}

\kamokumei{類題演習 解答例 --- 連立線形微分方程式}

\daimon{【解法】連立線形微分方程式の解き方と導出}

\noindent
2 元の連立 1 階線形微分方程式
\[
\frac{d}{dt}\begin{pmatrix} x \\ y \end{pmatrix}
= A \begin{pmatrix} x \\ y \end{pmatrix},
\qquad
A = \begin{pmatrix} a & b \\ c & d \end{pmatrix}
\]
（$A$ は定数行列）を考える。

\medskip
\noindent
\textbf{(1) 固有値・固有ベクトルによる解法}

\noindent
$A$ の固有値 $\lambda$ と固有ベクトル $\mathbf{v}$ を、特性方程式
\[
\det(A - \lambda I) = \lambda^{2} - (\operatorname{tr} A)\lambda + \det A = 0
\]
から求める。$A\mathbf{v} = \lambda \mathbf{v}$ を満たす $\mathbf{v} \neq \mathbf{0}$ に対して、
\[
\begin{pmatrix} x \\ y \end{pmatrix} = \mathbf{v}\, e^{\lambda t}
\]
は方程式を満たす（これは実際に代入すれば確かめられる：左辺 $= \lambda \mathbf{v} e^{\lambda t}$、右辺 $= A\mathbf{v} e^{\lambda t} = \lambda \mathbf{v} e^{\lambda t}$）。

\smallskip
\noindent
\emph{相異なる実固有値 $\lambda_{1} \neq \lambda_{2}$ の場合}: 対応する固有ベクトルを $\mathbf{v}_{1}, \mathbf{v}_{2}$ とすると、一般解は
\[
\begin{pmatrix} x \\ y \end{pmatrix}
= C_{1}\, \mathbf{v}_{1}\, e^{\lambda_{1} t} + C_{2}\, \mathbf{v}_{2}\, e^{\lambda_{2} t}
\]
（$\mathbf{v}_{1}, \mathbf{v}_{2}$ は一次独立なので、これが解空間の基底を与える）。

\smallskip
\noindent
\emph{複素固有値 $\lambda = \alpha \pm i\beta$ の場合}: 固有ベクトルも複素数となり $\mathbf{v} = \mathbf{p} + i\mathbf{q}$ とおける。複素解
\[
\mathbf{v}\,e^{(\alpha+i\beta)t}
= e^{\alpha t}(\mathbf{p} + i\mathbf{q})(\cos\beta t + i\sin\beta t)
\]
の実部・虚部はそれぞれ実数解となり、一般解は
\[
\begin{pmatrix} x \\ y \end{pmatrix}
= e^{\alpha t}\Big[ C_{1}(\mathbf{p}\cos\beta t - \mathbf{q}\sin\beta t) + C_{2}(\mathbf{p}\sin\beta t + \mathbf{q}\cos\beta t) \Big]
\]
（$C_1, C_2$ は実の任意定数）と表せる。

\medskip
\noindent
\textbf{(2) 消去法（1 変数の 2 階線形方程式への帰着）との対応}

\noindent
$x' = ax + by,\ y' = cx + dy$ （$b \neq 0$）から、第 1 式より $y = \dfrac{x' - ax}{b}$ とし、これを第 2 式に代入すると $x$ のみの 2 階線形方程式
\[
x'' - (a+d)x' + (ad - bc)x = 0
\]
が得られる。この特性方程式 $\mu^{2} - (a+d)\mu + (ad-bc) = 0$ は、行列 $A$ の特性方程式 $\lambda^{2} - (\operatorname{tr} A)\lambda + \det A = 0$ と完全に一致する（$\operatorname{tr} A = a+d$, $\det A = ad - bc$）。したがって消去法で得られる特性根は、行列の固有値そのものであり、両者は同一の解を与える。

\medskip
\noindent
初期条件が与えられた場合は、一般解に $t = 0$（または指定の初期時刻）を代入して定数 $C_{1}, C_{2}$ を定める。

\clearpage
\begin{mondaiboxS}{問1}
$\dfrac{dx}{dt} = x + 2y,\quad \dfrac{dy}{dt} = 2x + y$ の一般解を求めよ。
\end{mondaiboxS}
\kaitou
係数行列は $A = \begin{pmatrix} 1 & 2 \\ 2 & 1 \end{pmatrix}$。特性方程式は
\[
\det(A - \lambda I) = (1-\lambda)^{2} - 4 = \lambda^{2} - 2\lambda - 3 = (\lambda - 3)(\lambda + 1) = 0
\]
より固有値は $\lambda_{1} = 3,\ \lambda_{2} = -1$。

\noindent
$\lambda_{1} = 3$ に対して $(1-3)x + 2y = 0 \Rightarrow x = y$ より固有ベクトル $\mathbf{v}_{1} = \begin{pmatrix} 1 \\ 1 \end{pmatrix}$。

\noindent
$\lambda_{2} = -1$ に対して $(1+1)x + 2y = 0 \Rightarrow x = -y$ より固有ベクトル $\mathbf{v}_{2} = \begin{pmatrix} 1 \\ -1 \end{pmatrix}$。

\noindent
よって一般解は
\[
\boxed{\;
x(t) = C_{1} e^{3t} + C_{2} e^{-t}, \qquad
y(t) = C_{1} e^{3t} - C_{2} e^{-t}
\;}
\qquad (C_{1}, C_{2} \text{ は任意定数})
\]
（検算: $x' = 3C_{1}e^{3t} - C_{2}e^{-t}$ に対し $x + 2y = C_{1}e^{3t}+C_{2}e^{-t}+2C_{1}e^{3t}-2C_{2}e^{-t} = 3C_{1}e^{3t}-C_{2}e^{-t} = x'$。また $y' = 3C_{1}e^{3t}+C_{2}e^{-t}$ に対し $2x+y = 2C_{1}e^{3t}+2C_{2}e^{-t}+C_{1}e^{3t}-C_{2}e^{-t} = 3C_{1}e^{3t}+C_{2}e^{-t} = y'$。いずれも一致するので正しい。）

\clearpage
\begin{mondaiboxS}{問2}
$\dfrac{dx}{dt} = x - y,\quad \dfrac{dy}{dt} = x + y$ の一般解を求めよ。
\end{mondaiboxS}
\kaitou
係数行列は $A = \begin{pmatrix} 1 & -1 \\ 1 & 1 \end{pmatrix}$。特性方程式は
\[
\det(A - \lambda I) = (1-\lambda)^{2} + 1 = \lambda^{2} - 2\lambda + 2 = 0
\]
より固有値は $\lambda = 1 \pm i$。

\noindent
$\lambda = 1 + i$ に対して $(1 - (1+i))x - y = 0 \Rightarrow -ix - y = 0 \Rightarrow y = -ix$ より固有ベクトル $\mathbf{v} = \begin{pmatrix} 1 \\ -i \end{pmatrix}$。

\noindent
複素解 $\mathbf{v}\,e^{(1+i)t} = e^{t}(\cos t + i \sin t)\begin{pmatrix} 1 \\ -i \end{pmatrix}$ の実部・虚部を成分ごとに計算すると
\[
\begin{pmatrix} x \\ y \end{pmatrix}
= e^{t}\begin{pmatrix} \cos t \\ \sin t \end{pmatrix} + i\, e^{t}\begin{pmatrix} \sin t \\ -\cos t \end{pmatrix}
\]
したがって実部 $e^{t}(\cos t, \sin t)$ と虚部 $e^{t}(\sin t, -\cos t)$ はそれぞれ実数解であり、一般解は
\[
\boxed{\;
x(t) = e^{t}(C_{1}\cos t + C_{2}\sin t), \qquad
y(t) = e^{t}(C_{1}\sin t - C_{2}\cos t)
\;}
\qquad (C_{1}, C_{2} \text{ は任意定数})
\]
（検算: $x' = e^{t}(C_{1}\cos t+C_{2}\sin t) + e^{t}(-C_{1}\sin t+C_{2}\cos t) = e^{t}[(C_{1}+C_{2})\cos t+(C_{2}-C_{1})\sin t]$。一方 $x-y = e^{t}(C_{1}\cos t+C_{2}\sin t) - e^{t}(C_{1}\sin t-C_{2}\cos t) = e^{t}[(C_{1}+C_{2})\cos t+(C_{2}-C_{1})\sin t]$ で一致。$y' = e^{t}(C_{1}\sin t-C_{2}\cos t)+e^{t}(C_{1}\cos t+C_{2}\sin t) = e^{t}[(C_{1}-C_{2})\cos t+(C_{1}+C_{2})\sin t]$。一方 $x+y = e^{t}(C_{1}\cos t+C_{2}\sin t)+e^{t}(C_{1}\sin t-C_{2}\cos t) = e^{t}[(C_{1}-C_{2})\cos t+(C_{1}+C_{2})\sin t]$ で一致。よって正しい。）

\clearpage
\begin{mondaiboxS}{問3}
$\dfrac{dx}{dt} = 4x - 2y,\quad \dfrac{dy}{dt} = x + y,\quad x(0) = 1,\ y(0) = 0$ を解け。
\end{mondaiboxS}
\kaitou
係数行列は $A = \begin{pmatrix} 4 & -2 \\ 1 & 1 \end{pmatrix}$。特性方程式は
\[
\det(A - \lambda I) = (4-\lambda)(1-\lambda) + 2 = \lambda^{2} - 5\lambda + 6 = (\lambda-2)(\lambda-3) = 0
\]
より固有値は $\lambda_{1} = 3,\ \lambda_{2} = 2$。

\noindent
$\lambda_{1} = 3$ に対して $(4-3)x - 2y = 0 \Rightarrow x = 2y$ より固有ベクトル $\mathbf{v}_{1} = \begin{pmatrix} 2 \\ 1 \end{pmatrix}$。

\noindent
$\lambda_{2} = 2$ に対して $(4-2)x - 2y = 0 \Rightarrow x = y$ より固有ベクトル $\mathbf{v}_{2} = \begin{pmatrix} 1 \\ 1 \end{pmatrix}$。

\noindent
一般解は
\[
x(t) = 2C_{1}e^{3t} + C_{2}e^{2t}, \qquad
y(t) = C_{1}e^{3t} + C_{2}e^{2t}
\]
初期条件 $x(0) = 1,\ y(0) = 0$ を代入すると
\[
2C_{1} + C_{2} = 1, \qquad C_{1} + C_{2} = 0
\]
これを解いて $C_{1} = 1,\ C_{2} = -1$。したがって
\[
\boxed{\;
x(t) = 2e^{3t} - e^{2t}, \qquad
y(t) = e^{3t} - e^{2t}
\;}
\]
（検算: $x(0) = 2-1 = 1$, $y(0) = 1-1 = 0$ で初期条件を満たす。$x' = 6e^{3t}-2e^{2t}$ に対し $4x-2y = 4(2e^{3t}-e^{2t})-2(e^{3t}-e^{2t}) = 8e^{3t}-4e^{2t}-2e^{3t}+2e^{2t} = 6e^{3t}-2e^{2t} = x'$。$y' = 3e^{3t}-2e^{2t}$ に対し $x+y = (2e^{3t}-e^{2t})+(e^{3t}-e^{2t}) = 3e^{3t}-2e^{2t} = y'$。いずれも一致するので正しい。）

\end{document}
